%! Polynomial Functions and Complex Zeros — Unit 2, Lesson 2.4 — Guided Notes
\documentclass[10pt]{article}
\usepackage{precalculus-article}
\usepackage{precalculus-boxes}
\newcommand{\TallMath}[1]{$\displaystyle #1\rule[-1.4em]{0pt}{3.2em}$}

\begin{document}

\pageheader{Unit 2, Lesson 2.4}{Guided Notes}
\namedateperiod

\vspace{0.1in}

\begin{objectivebox}
Students will identify and interpret zeros of polynomial functions (real and complex),
including multiplicity, and determine whether a polynomial function is even or odd.
\begin{enumerate}[leftmargin=1.5em, itemsep=2pt]
  \item[\textbullet] \writeline
  \item[\textbullet] \writeline
  \item[\textbullet] \writeline
\end{enumerate}
\end{objectivebox}

\begin{vocabbox}
\termblanklong{zero (root)}
\termblanklong{multiplicity}
\termblanklong{complex zero}
\termblanklong{complex conjugate pair}
\termblanklong{even function}
\termblanklong{odd function}
\end{vocabbox}

\begin{hookbox}
Consider $p(x)=(x-2)^2(x+1)$.
\begin{itemize}[leftmargin=1.2em, itemsep=3pt]
  \item Where does the graph of $y=p(x)$ \emph{cross} the $x$-axis? \blank{4cm}
  \item Where does it \emph{touch and bounce} without crossing? \blank{4cm}
  \item What is it about the \emph{exponent} on a factor that tells you which will happen?
        \writeline
\end{itemize}
\end{hookbox}

\begin{notesbox}{Zeros and Linear Factors}
\textbf{Key connection:} For a polynomial $p(x)$, if $p(a) = \blank{1.5cm}$, then $a$
is called a \blank{3cm} of $p$, and $(x - a)$ is a \blank{3cm} factor of $p$.

\vspace{0.15in}

\textbf{Multiplicity:} If the factor $(x-a)$ appears \blank{2cm} times, then $a$ has
multiplicity \blank{1.5cm}. A polynomial of degree $n$ has exactly \blank{1.5cm}
complex zeros (counting multiplicities).

\vspace{0.1in}

\renewcommand{\arraystretch}{1.6}
\begin{tabularx}{\linewidth}{|X|X|X|X|}
\hline
\rowcolor{sky}
\textbf{Polynomial} & \textbf{Factored Form} & \textbf{Zeros} & \textbf{Multiplicities} \\
\hline
$p(x)=(x-2)^2(x+1)$ & & & \\[14pt]
\hline
$q(x)=x^3-x^2-6x$ & & & \\[14pt]
\hline
$r(x)=(x+1)^3(x-4)^2$ & already factored & & \\[14pt]
\hline
\end{tabularx}
\end{notesbox}

\begin{notesbox}{Graph Behavior at Real Zeros}
\renewcommand{\arraystretch}{1.6}
\begin{tabularx}{\linewidth}{|l|X|X|}
\hline
\rowcolor{sky}
\textbf{Multiplicity} & \textbf{Graph behavior at the zero} & \textbf{Example} \\
\hline
Odd (1, 3, 5, \ldots) & \blank{4cm} the $x$-axis & $(x+1)^1$ in $p(x)$ above \\
\hline
Even (2, 4, 6, \ldots) & \blank{4cm} the $x$-axis (tangent) & $(x-2)^2$ in $p(x)$ above \\
\hline
\end{tabularx}

\vspace{0.1in}
For $p(x)=(x+2)(x-1)^2(x-3)$, predict the graph's behavior at each zero:

\begin{tabularx}{\linewidth}{|X|X|X|}
\hline
\rowcolor{sky}
\textbf{Zero} & \textbf{Multiplicity} & \textbf{Crosses or tangent?} \\
\hline
$x = -2$ & \blank{2cm} & \blank{3cm} \\
\hline
$x = 1$  & \blank{2cm} & \blank{3cm} \\
\hline
$x = 3$  & \blank{2cm} & \blank{3cm} \\
\hline
\end{tabularx}
\end{notesbox}

\begin{notesbox}{Complex Zeros \& Conjugate Pairs}
\textbf{Complex Conjugate Pairs Theorem:} If $p(x)$ is a polynomial with
\blank{3cm} coefficients and $a+bi$ (with $b\neq 0$) is a zero, then
\blank{3cm} is also a zero of $p$.

\vspace{0.1in}
\textbf{Example:} Suppose $p(x)$ has real coefficients and $2+3i$ is a zero.
\begin{itemize}[leftmargin=1.2em, itemsep=4pt]
  \item The other complex zero must be: \blank{4cm}
  \item Write the quadratic factor formed by this conjugate pair:\\[4pt]
        $(x-(2+3i))(x-(2-3i)) = x^2 - \blank{2cm}x + \blank{2cm}$
  \item If $p$ also has real zero $x=1$, what is the minimum degree of $p$?
        \blank{2cm}
\end{itemize}
\end{notesbox}

\begin{notesbox}{Even and Odd Functions}
\textbf{Definitions:}
\begin{itemize}[leftmargin=1.2em, itemsep=4pt]
  \item $f$ is \textbf{even} if $f(-x) = \blank{2.5cm}$ for all $x$ in its domain.
        Graph is symmetric about the \blank{3cm}.
  \item $f$ is \textbf{odd} if $f(-x) = \blank{2.5cm}$ for all $x$ in its domain.
        Graph is symmetric about the \blank{3cm}.
\end{itemize}

\vspace{0.1in}
\textbf{Algebraic test — steps:}
\begin{enumerate}[leftmargin=1.5em, itemsep=3pt]
  \item Replace every $x$ with \blank{1.5cm} and simplify.
  \item Compare the result to \blank{3cm} (even) or \blank{3cm} (odd).
  \item If neither equality holds, the function is \blank{2.5cm}.
\end{enumerate}

\vspace{0.1in}
\renewcommand{\arraystretch}{1.6}
\begin{tabularx}{\linewidth}{|X|X|X|}
\hline
\rowcolor{sky}
\textbf{Function} & \textbf{$f(-x) = \ldots$} & \textbf{Even, Odd, or Neither?} \\
\hline
$f(x)=x^4-2x^2$ & & \\[12pt]
\hline
$g(x)=x^3-x$ & & \\[12pt]
\hline
$h(x)=x^2+x$ & & \\[12pt]
\hline
\end{tabularx}
\end{notesbox}

\begin{practicebox}
\begin{enumerate}[leftmargin=1.5em, itemsep=10pt]
  \item A polynomial with real coefficients has zeros $x=4$ (multiplicity 2) and $x=1-i$.
        \begin{enumerate}[leftmargin=1.5em, itemsep=4pt]
          \item List all zeros of this polynomial.
                \writeline
          \item What is the minimum degree?  \blank{2cm}
          \item At which real zero does the graph touch but not cross the $x$-axis?
                \blank{3cm}
        \end{enumerate}

  \item Determine whether $k(x)=2x^6-3x^4+x^2$ is even, odd, or neither.
        Show the algebraic test.
        \writelines{3}
\end{enumerate}
\end{practicebox}

\end{document}
